\section{Introducción}
El monitoreo automatizado de la fauna silvestre mediante imágenes aéreas es crucial para la conservación ecológica, pero enfrenta un desafío crítico en la sabana africana: la detección precisa en escenarios de alta densidad. Especies gregarias como elefantes y búfalos forman manadas compactas con severa oclusión y superposición entre individuos, condiciones donde los sistemas de visión por computador convencionales fallan sistemáticamente.

El problema fundamental radica en las arquitecturas tradicionales basadas en CNN y anclas, que dependen del algoritmo de Supresión de No Máximos (NMS). En manadas densas, el NMS elimina erróneamente detecciones válidas de animales contiguos, provocando un subconteo crónico. Aunque alternativas como HerdNet \cite{herdnet} evitan el NMS mediante la estimación de densidad por puntos, su dependencia de backbones CNN limita la captura del contexto global necesario para una clasificación robusta, afectando especialmente a las especies minoritarias.

Este trabajo se justifica por la necesidad de una solución que supere simultáneamente las limitaciones de oclusión y discriminación semántica. El objetivo principal es implementar y evaluar la arquitectura \textbf{RF-DETR} (\textit{Real-Time Detection Transformer}) \cite{rf-detr} en este contexto desafiante. Nuestra hipótesis es que el mecanismo de atención de los Transformers, capaz de modelar dependencias globales, junto con la predicción de conjuntos \textit{end-to-end} que elimina por completo la necesidad del NMS, ofrece una solución estructural al problema del subconteo en manadas.

Las principales contribuciones de este estudio son:  (1) Una comparación de tres variantes de RF-DETR (Nano, Small, Large) frente al baseline HerdNet. (2) Un análisis del compromiso entre precisión en alta densidad y costo computacional de las arquitecturas, identificando configuraciones óptimas para el despliegue práctico. (3) La validación de una estrategia de entrenamiento de dos fases con Minería de Negativos Difíciles (HNM) para refinar la precisión en fondos complejos.
